The eight-factor ontology is an adapted compression of recurring cause families in Assaf and Al-Hejji, Sambasivan and Soon, and Doloi \textit{et al.}, mapped onto NYC administrative reason language rather than survey Likert items~\cite{Assaf2006,Sambasivan2007,Doloi2012}. Factor presence for a PID is the logical OR of mapped reason tokens across that PID's history rows retained after blank-reason exclusion. This presence definition is intentionally coarse: it matches the binary parent nodes of the Noisy-OR child and avoids fabricating graded severities that the open-data reasons do not support. Parameters $(p_c,\ell)=(0.55,0.02)$ are TRAIN-fit values used in the released BN configuration; they are not transferred from clinical Noisy-OR tables~\cite{Onisko2001}. Sensitivity of early-warning accuracy to prefix length is reported in Results; BN cause-probability sensitivity is secondary because policy ties persist under the clamped evidence model even when posterior levels shift slightly.

The early-warning champion is a histogram-based gradient boosting classifier trained on prefix features. Candidate pipelines evaluated on VAL included the HGB model and a stacking ensemble; the HGB pipeline was retained after VAL accuracy/recall screening. Hyperparameters follow scikit-learn defaults unless overridden in the training script, with the binding selection step being threshold $\tau$ on VAL rather than exhaustive test-set tuning. This protocol intentionally understates ``fully optimized champion versus neglected baseline'' risk by keeping the primary comparison against published literature accuracy references and against BN policies on the same TEST PIDs~\cite{Egwim2021,Friedman2001}.

Budgeted evaluation replays each TEST PID's inspectable-factor set under four policies for $B\in\{1,2,3,4\}$. For each policy and budget, the system records posterior trajectories, alert decisions at $\tau$, and summary metrics (recall, precision, FPR, ROC-AUC, Brier, mean steps to first alert among alerting severe cases). Expected utility and realized alert loss are computed from the same asymmetric costs used inside VoI, ensuring that the decision criterion and the evaluation criterion are aligned~\cite{Howard1966,Vickers2006}. Randomness in bootstrap intervals uses a fixed seed schedule with $400$ resamples of TEST PIDs; classifier training uses a fixed TRAIN/VAL/TEST PID partition rather than repeated random splits, because PID-held-out integrity is the higher-priority design constraint for this study~\cite{Tagharobi2025}.

Reason$\rightarrow$factor mapping quality is measured on VAL texts. Regex rules fire first when high-confidence patterns match; otherwise a TF--IDF logistic model trained only on TRAIN weak labels provides a fallback. Macro-F1 of the NLP-only path is modest because unique reason strings are few and class support is uneven; the operational regex-first mapper recovers high weighted F1 on rows where regex is informative. This dual-path design is part of the method, not a post-hoc excuse: the BN parents must be populated by an auditable extractor before VoI can act~\cite{Onisko2001,Woldesenbet2016}.

Fig.~\ref{fig:framework} is the sole full-width architecture figure. Ordinary quantitative figures in Results are generated from saved prediction tables, bootstrap arrays, calibration bins, decision-curve points, and policy scoreboards produced by the experiment scripts. No quantitative plot is drawn from hypothesized numbers. The experimental evidence trail is therefore: NYC Open Data endpoints $\rightarrow$ PID split and feature tables $\rightarrow$ trained models and BN replay logs $\rightarrow$ aggregated metrics $\rightarrow$ figures/tables $\rightarrow$ claims in Section~\ref{sec:results}.
